% SPDX-License-Identifier: GPL-3.0-or-later
% Copyright (C) 2026 Alessandro Greco
\chapter{Tor Browser, Tails и Whonix}\label{cap:23-browser}
\aggiornatoal{23 сентября 2026 г.}{документ по архитектуре Tor Browser (2018); руководство Tor Browser; документация Tails и Whonix}

\begin{obiettivi}
К концу этой главы вы сможете:
\begin{itemize}[nosep]
  \item объяснять, почему недостаточно просто настроить обычный браузер на использование Tor;
  \item перечислять требования безопасности и приватности, лежащие в основе архитектуры Tor Browser;
  \item связывать изоляцию по сайтам в Tor Browser с изоляцией потоков tor;
  \item сравнивать, что защищают Tor Browser, Tails и Whonix и какой ценой достигается каждая модель.
\end{itemize}
\end{obiettivi}

Модель угроз 2004 года прямо указывала, что Tor не нормализует прикладные протоколы (\cref{sec:minaccia-nonobiettivi}). Если браузер раскрывает вашу идентичность, луковая маршрутизация не сможет вас защитить. В этой главе рассматриваются инструменты, созданные именно для этого уровня: Tor Browser, который согласует поведение браузера с защитными свойствами сети, и две операционные системы — Tails и Whonix, распространяющие этот подход на весь компьютер.

% ---------------------------------------------------------------------------
\section{Почему одного прокси недостаточно}\label{sec:browser-proxy}

Настройка обычного браузера на использование SocksPort tor во многих случаях скрывает IP-адрес, но не во всех и почти ничего не скрывает сверх этого:
\begin{itemize}
  \item Некоторые компоненты могут подключаться в обход прокси: локальное DNS-разрешение, WebRTC, расширения или внешние программы, запускаемые браузером.
  \item Файлы cookie и другие данные состояния браузера связывают сегодняшние посещения со вчерашними, возможно выполненными без Tor.
  \item Сайт может собирать десятки характеристик браузера — размер окна, установленные шрифты, язык, часовой пояс, особенности отрисовки — и объединять их в \textbf{цифровой отпечаток}, отличающий одного пользователя от миллионов.
  \item Если все вкладки используют одну цепь, отслеживающий компонент, встроенный в несколько сайтов, может наблюдать запросы, приходящие одновременно с одного выходного реле.
\end{itemize}
Первая проблема угрожает адресу, остальные — \emph{несвязываемости} действий, одной из центральных целей Tor.

% ---------------------------------------------------------------------------
\section{Tor Browser}\label{sec:browser-tb}

Tor Browser — модифицированная версия Firefox Extended Support Release от Mozilla, распространяемая Tor Project вместе с tor. Его архитектура описана в документе, который уже не актуален во всех деталях, но остаётся важным изложением основных принципов~\cite{perry2018torbrowser}. Требования делятся на две группы.

\paragraph{Требования безопасности.} Они определяют безопасное использование Tor; их нарушение может немедленно раскрыть пользователя:
\begin{description}
  \item[Обязательное использование прокси] Никакой контент не должен обходить tor.
  \item[Разделение состояния] Информация не должна совместно использоваться с другими браузерами или сеансами вне Tor.
  \item[Отсутствие постоянных следов] Ничто, раскрывающее историю просмотра, не должно оставаться на диске или в памяти после завершения сеанса, если пользователь явно не решил сохранить это.
  \item[Изоляция приложения] Всё, что браузер оставляет в системе, должно удаляться вместе с ним.
\end{description}

\paragraph{Требования приватности.} Они не позволяют связывать разные действия и организованы вокруг \textbf{сайта в адресной строке (источника, origin)}, то есть сайта, который посещает пользователь:
\begin{description}
  \item[Идентификаторы] Данные, сохранённые третьей стороной, например отслеживающим компонентом, встроенным в несколько сайтов, во время посещения одного сайта не должны становиться доступными той же стороне при её появлении на другом.
  \item[Цифровые отпечатки] Характеристики браузера не должны позволять связывать действия на разных сайтах.
  \item[Долгосрочная несвязываемость] Пользователь должен иметь возможность одной простой командой удалить всё состояние и получить новую идентичность.
\end{description}
Документ предлагает и полезный принцип проектирования: защитные механизмы не должны нарушать ожидаемое пользователем поведение браузера. Предшественник Tor Browser, расширение Torbutton, позволял включать и выключать Tor в одном браузере. Пользователи путали вкладки, открытые в разных режимах. Интерфейс, провоцирующий ошибки, сам становится слабостью безопасности.

\subsection{Одна цепь на сайт}

Изоляция по источнику сайта распространяется и на сеть. Tor Browser передаёт идентификатор сайта из адресной строки в SOCKS-учётных данных, отправляемых tor. По умолчанию tor помещает потоки с разными учётными данными (\opt{IsolateSOCKSAuth}) в разные цепи. Этот механизм мы проверили в Лабораторной работе~11 (\cref{sec:lab-stem}). С точки зрения пользователя, как описано в руководстве~\cite{tbmanual}, два сайта, содержащие один и тот же отслеживающий компонент, используют разные цепи, тогда как вкладки одного сайта разделяют одну.

На эту схему действуют две команды:
\begin{description}
  \item[Новая цепь для этого сайта (New Circuit for This Site)] Перезагружает текущий сайт через новую цепь, например когда выходное реле не может к нему подключиться. Состояние браузера не очищается, а цепи других сайтов не меняются.
  \item[Новая личность (New Identity)] Закрывает все вкладки, удаляет файлы cookie, историю и другое состояние и получает новые цепи для всей активности, реализуя требование долгосрочной несвязываемости.
\end{description}
Ни одно из этих действий само по себе не меняет входное реле (guard), который остаётся стабильным месяцами по причинам, рассмотренным в \cref{sec:guard}.

Изоляция по сайтам означает прежде всего \emph{разделение разных сайтов}, а не гарантию единственной цепи для сайта на весь сеанс: цепи со временем заменяются. Команда \enquote{Новая цепь для этого сайта} перезагружает активную вкладку; другие вкладки того же сайта перейдут на новый маршрут после перезагрузки. Текущие соединения не обязательно разрываются. \enquote{Новая личность} также прекращает текущую активность и загрузки. Ни одна из команд не предназначена для принудительной смены входного реле: tor управляет им отдельно.

\subsection{Защита от цифровых отпечатков}

Стратегия Tor Browser против снятия цифровых отпечатков состоит не в том, чтобы сделать каждого пользователя уникальным, а в том, чтобы пользователи выглядели \emph{одинаково}. Если все пользователи Tor Browser сообщают одинаковые значения, отпечаток говорит немногим больше, чем \enquote{это пользователь Tor Browser}. Отсюда следуют несколько заметных решений. Окна открываются с округлёнными размерами, а при изменении размера вокруг содержимого страницы добавляются поля, чтобы сообщаемые размеры оставались стандартизированными кратными значениями (\emph{letterboxing — добавление полей}). Сообщаемый часовой пояс — UTC. Многие интерфейсы, раскрывающие детали операционной системы, отключены или стандартизированы. Этим же объясняется рекомендация не устанавливать расширения и не менять настройки: каждое отличие уменьшает множество пользователей, среди которых можно скрыться, то есть набор анонимности, обсуждавшийся в \cref{cap:01-anonimato}.

Цель защиты — \emph{уменьшить число наблюдаемых отличий}, а не сделать все отпечатки буквально одинаковыми. Платформа, язык, уровень безопасности и другие свойства могут различать группы пользователей. Поэтому речь идёт о снижении уникальности, а не о гарантии полного совпадения.

\subsection{Уровни безопасности}

Многие уязвимости браузера используют JavaScript и сложные функции. Tor Browser предлагает три уровня~\cite{tbmanual}:
\begin{description}
  \item[Standard] Все функции включены; уровень по умолчанию.
  \item[Safer] JavaScript отключён на сайтах без HTTPS; некоторые шрифты и математические символы отключены; аудио и видео запускаются только после клика.
  \item[Safest] JavaScript отключён везде, а изображения и другой контент также могут быть ограничены. Статические сайты остаются доступными.
\end{description}
Выбор уровня — компромисс между безопасностью и удобством. Сам выбранный уровень может влиять на цифровой отпечаток, поскольку сайт способен косвенно определить его.

\begin{approfondimento}[Воспроизводимые сборки]
Пользователям, загружающим Tor Browser, нужна уверенность, что распространяемый бинарный файл соответствует опубликованному исходному коду. Tor Project создаёт \emph{воспроизводимые сборки}: одинаковые исходники и среда сборки должны давать побитно идентичные файлы. Документ по архитектуре описывает работу, необходимую для устранения источников различий, включая даты, пути файлов и порядок элементов архивов~\cite{perry2018torbrowser}. Независимая повторная сборка и сравнение результатов позволяют обнаружить вредоносные изменения, внесённые в среду сборки дистрибутива.
\end{approfondimento}

% ---------------------------------------------------------------------------
\section{За пределами браузера}\label{sec:browser-sistemi}

Tor Browser защищает собственный трафик, но не трафик других приложений на том же компьютере. Эти программы всё ещё могут подключаться напрямую, а ошибка или эксплуатация уязвимости браузера может раскрыть реальный адрес. Две операционные системы распространяют защиту на всю машину с помощью разных архитектур (\cref{fig:architetture}).

\begin{figure}[htbp]
\centering
% SPDX-License-Identifier: GPL-3.0-or-later
% Copyright (C) 2026 Alessandro Greco
%% Русское альтернативное описание EPUB.
\altfig{Сравнение трёх архитектур. (a) Tor Browser направляет через tor только собственный трафик; другие приложения на компьютере не защищены автоматически и могут подключаться напрямую. Пунктирная рамка охватывает только браузер и tor. (b) Tails загружается как live-система с USB и работает преимущественно в памяти; правила межсетевого экрана обычно блокируют прямые соединения приложений и заставляют использовать Tor. Существуют намеренные исключения, включая небезопасный браузер Unsafe Browser для captive portal, который не использует Tor. (c) Whonix разделяет приложения в Workstation VM и tor в Gateway VM, связанные изолированной виртуальной сетью. Компрометация хостовой системы или гипервизора по-прежнему опасна.}
% Схема: браузер, live-система с межсетевым экраном, две VM
\begin{tikzpicture}[font=\sffamily\scriptsize,
  box/.style={draw,rounded corners=2pt,align=center,minimum width=2.2cm,minimum height=0.7cm},
  conf/.style={draw=torviola,dashed,rounded corners=4pt,inner sep=5pt}]
  % Tor Browser
  \node[box] (tb) at (0,1.6) {Tor Browser};
  \node[box,fill=torviola!15] (t1) at (0,0.4) {tor};
  \node[box,fill=torgrigio!15] (a1) at (0,-0.8) {другие приложения};
  \draw[->] (tb) -- node[right] {SOCKS} (t1);
  \draw[->,avviso] (a1.east) -- ++(0.7,0) node[right,align=left] {возможен прямой\\доступ в сеть};
  \node[conf,fit=(tb)(t1)] (s1) {};
  \node[above=2pt of s1] {\small (a) только трафик браузера};
  % Tails
  \begin{scope}[xshift=4.4cm]
  \node[box] (tb2) at (0,1.6) {приложения};
  \node[box,fill=torviola!15] (t2) at (0,-0.8) {tor};
  \node[box,fill=avviso!10] (fw) at (0,0.4) {firewall: блокирует обход};
  \draw[->] (tb2) -- (fw);
  \draw[->] (fw) -- (t2);
  \node[conf,fit=(tb2)(fw)] (s2) {};
  \node[above=2pt of s2] {\small (b) Tails: live USB, RAM};
  \end{scope}
  % Whonix
  \begin{scope}[xshift=8.8cm]
  \node[box] (ws) at (0,1.6) {Workstation\\(приложения)};
  \node[box,fill=torviola!15] (gw) at (0,-0.8) {Gateway\\(tor)};
  \draw[->] (ws) -- node[right,align=left] {изолированная\\виртуальная сеть} (gw);
  \node[conf,fit=(ws)] {};
  \node[conf,fit=(gw)] {};
  \node[above=12pt of ws] {\small (c) Whonix: две VM};
  \end{scope}
  \foreach \n in {t1,t2,gw} \draw[->,thick,torverde!60!black] (\n.south) -- ++(0,-0.5) node[below] {сеть Tor};
\end{tikzpicture}

\caption{Три модели маршрутизации и изоляции, не абсолютные гарантии анонимности. (a) Tor Browser пропускает собственный трафик через tor, но другие программы на хосте могут подключаться напрямую. (b) Tails обычно блокирует обход Tor межсетевым экраном; \emph{Unsafe Browser} для captive portal — явное неанонимное исключение. Tails загружается с USB как live-система, работает преимущественно в RAM и поддерживает необязательное постоянное хранилище. (c) Whonix размещает Workstation с приложениями и Gateway с tor в двух VM, связанных изолированной виртуальной сетью; компрометация хоста или гипервизора может нарушить эту защиту.}
\label{fig:architetture}
\end{figure}

\subsection{Tails}

Tails — основанная на Debian операционная система, загружаемая с USB-накопителя~\cite{tailsabout}. У неё два определяющих свойства:
\begin{description}
  \item[Всё через Tor] Приложения, пытающиеся выйти в Интернет без tor, блокируются.
  \item[Амнезия] Система работает в памяти, не записывает данные на диск компьютера и забывает своё состояние после выключения. Данные, которые пользователь решил сохранить, помещаются в необязательное зашифрованное \emph{постоянное хранилище} на USB-накопителе.
\end{description}
Tails включает Tor Browser и другие осторожно настроенные приложения. Проект присоединился к Tor Project в сентябре 2024 года~\cite{torblogtails}. Документация прямо указывает ограничения: Tails не защищает от модифицированного аппаратного обеспечения, например аппаратного регистратора нажатий клавиш, или от USB-накопителя, подготовленного на уже заражённом компьютере.

Эти свойства относятся к обычному режиму работы Tails и \emph{не дают абсолютной гарантии} отсутствия следов на любых носителях. Пользователь может самостоятельно подключить диски или создать постоянные файлы; вредоносное ПО, скомпрометированное оборудование и действия вне Tails тоже выходят за границы модели угроз.

\subsection{Whonix}

Whonix разделяет функции между двумя виртуальными машинами. \textbf{Gateway} запускает tor и является единственной виртуальной машиной, подключённой к внешней сети. \textbf{Workstation} запускает приложения и соединяется только с Gateway через внутреннюю виртуальную сеть. Поэтому весь трафик Workstation обязан проходить через tor благодаря архитектуре, а не только настройкам приложений. По сравнению с Tails это даёт более сильное разделение: даже атакующий, получивший административный контроль над Workstation, не может напрямую узнать реальный адрес, поскольку сама машина его не знает. Цена — зависимость от безопасности ПО виртуализации и отсутствие амнезии, если Whonix не запускается внутри системы, которая её обеспечивает.

Разделение шлюза и рабочей станции — важная граница безопасности, но \emph{не абсолютная гарантия} при компрометации рабочей станции. Атакующий может раскрыть личность через данные приложений; взлом хоста, гипервизора или шлюза способен обойти сетевую изоляцию. Сохраняют значение корректная конфигурация и защита межсетевого экрана.

\paragraph{Фильтрация порта управления.} Некоторые приложения Workstation, включая функцию New Identity в Tor Browser, должны взаимодействовать с ControlPort tor на шлюзе (Gateway). Однако протокол управления (\cref{sec:stem-protocollo}) также позволяет запросить публичный адрес tor через \texttt{GETINFO address} — именно то, чего Workstation знать не должна. Поэтому Whonix помещает между ними \emph{onion-grater}, прокси, разрешающий только явно допустимые команды~\cite{whonixgrater}. Функция, безопасная в одном контексте, может стать утечкой в другом.

\begin{attenzione}[Инструменты не заменяют осторожное поведение]
Tor Browser, Tails и Whonix уменьшают число возможных технических ошибок, но не устраняют поведенческие: вход в личный аккаунт, открытие загруженных документов вне защищённой среды или раскрытие идентифицирующих сведений. Этот уровень рассматривается в \cref{cap:28-opsec}.
\end{attenzione}

% ---------------------------------------------------------------------------
\section*{Итоги}
\begin{itemize}
  \item Обычный браузер за Tor может пропускать трафик мимо прокси и, что ещё важнее, связывать действия через общее состояние, цифровые отпечатки или цепи.
  \item Tor Browser следует явным требованиям архитектуры: обязательное использование прокси, разделение состояния, ограниченное долговременное хранение данных, несвязываемость между сайтами и доступная функция смены личности (New Identity).
  \item Изоляция по сайтам основана на SOCKS-учётных данных и механизме изоляции потоков tor.
  \item Tor Browser противодействует цифровым отпечаткам, стандартизируя наблюдаемые характеристики пользователей; уровни безопасности уменьшают поверхность атаки.
  \item Tails принудительно направляет весь трафик через Tor и избегает постоянных следов; Whonix изолирует приложения внутри машины, не знающей реального сетевого адреса.
\end{itemize}

% ---------------------------------------------------------------------------
\section*{Упражнения}

\begin{esercizio}[1]\label{es:23-identita}
С точки зрения цепей и состояния браузера, чем отличается \enquote{Новая цепь для этого сайта} от \enquote{Новая личность}? Почему ни одна из команд не меняет входное реле (guard)?
\end{esercizio}

\begin{esercizio}[1]\label{es:23-estensioni}
Пользователь устанавливает в Tor Browser расширение для блокировки рекламы, считая, что это улучшит приватность. Что он может выиграть и что потерять?
\end{esercizio}

\textbf{Методическое замечание.} Лабораторная работа~11 имитирует разные учётные данные SOCKS, но не открывает три настоящие вкладки Tor Browser. Совпадающие учётные данные могут привести к повторному использованию цепи, однако это не гарантировано постоянно: цепи устаревают, возможны дополнительные правила изоляции.

\begin{esercizio}[2]\label{es:23-isolamento}
Используя лабораторную сеть и скрипт Лабораторной работы~11, смоделируйте три вкладки: две посещают один сайт, третья — другой, причём имя сайта используется как учётные данные SOCKS. Через порт управления (ControlPort) проверьте, какие потоки используют одну цепь.
\end{esercizio}

\begin{esercizio}[2]\label{es:23-scelta}
Для каждого сценария выберите подходящий инструмент из Tor Browser, Tails и Whonix и объясните решение: (a) человек просматривает заблокированные сайты с домашнего компьютера; (b) журналист обрабатывает конфиденциальные документы на чужих компьютерах; (c) исследователь запускает несколько долгоживущих программ, которые все должны использовать Tor.
\end{esercizio}

\begin{esercizio}[3]\label{es:23-impronta}
Выберите пять свойств браузера, доступных сайту через JavaScript, и оцените, сколько бит информации каждое раскрывает в обычном браузере. Как Tor Browser стандартизирует эти свойства и какие различия между пользователями всё ещё остаются?
\end{esercizio}

% ---------------------------------------------------------------------------
\section*{Дополнительная литература}
\begin{itemize}
  \item M.~Perry, E.~Clark, S.~Murdoch, G.~Koppen, \titolo{The Design and Implementation of the Tor Browser}~\cite{perry2018torbrowser}.
  \item Руководство Tor Browser~\cite{tbmanual}.
  \item Документация Tails~\cite{tailsabout} и Whonix, особенно страница onion-grater~\cite{whonixgrater}.
\end{itemize}
